% experiments.tex -- Experiments

We applied stochastic attention to protein sequence generation across seven Pfam families and conducted a systematic scaling study to characterize performance as a function of family size. All experiments used the unadjusted Langevin algorithm (ULA) with step size $\alpha=0.01$. For each family and condition, we ran 30 independent chains of $T=5{,}000$ iterations, initialized near randomly chosen stored patterns with Gaussian perturbation ($\sigma_{\mathrm{init}}=0.01$). After discarding $2{,}000$ burn-in iterations and thinning every 100 steps, we retained 5 evenly spaced samples per chain for a total of $S=150$ generated sequences per condition. MALA chains were run in parallel with identical parameters to provide acceptance-rate diagnostics. Throughout, we used PCA dimensionality reduction retaining 95\% of the variance in the one-hot encoded alignment, followed by unit-norm projection to form the memory matrix $\hat{\mathbf{X}}$.

The inverse temperature was set automatically for each family via the entropy inflection criterion (Section~\ref{sec:beta-star}). We identified $\beta^*$ from the attention entropy curve and operated in two regimes: a \emph{generation} regime at $\beta_{\mathrm{gen}} \approx 2\beta^*$, where the sampler explores broadly while respecting family statistics, and a \emph{retrieval} regime at $\beta_{\mathrm{ret}} \approx 20\beta^*$, where samples remain close to stored patterns. Three simple baselines were evaluated under identical conditions: bootstrap replay (uniform resampling of stored patterns), Gaussian perturbation (stored pattern plus isotropic noise matched to the SA noise scale), and random convex combination (Dirichlet-weighted average of all stored patterns).

We evaluated generated sequences using five metrics. \emph{Novelty} measured departure from stored patterns as $1 - \max_k \cos(\hat{\boldsymbol{\xi}}, \mathbf{m}_k)$ in PCA space. \emph{Diversity} was the mean pairwise cosine distance among generated samples. \emph{Sequence identity} was the maximum fraction of matching residues between a generated sequence and any stored sequence, computed after argmax decoding from PCA space back to amino acid sequences. \emph{KL divergence} of amino acid composition quantified how well generated sequences preserved the family's residue frequency distribution. \emph{Valid residue fraction} confirmed that all decoded positions were standard amino acids.

\paragraph{Multiple Pfam families.}
We selected seven families spanning a range of sizes, sequence lengths, and conservation levels (Table~\ref{tab:families}): RRM (PF00076, $K{=}68$, $L{=}71$), SH3 (PF00018, $K{=}55$, $L{=}48$), WW (PF00397, $K{=}420$, $L{=}31$), Kunitz (PF00014, $K{=}99$, $L{=}53$), zinc finger C2H2 (PF00096, $K{=}151$, $L{=}23$), PDZ (PF00595, $K{=}44$, $L{=}83$), and protein kinase (PF00069, $K{=}37$, $L{=}262$). Seed alignments were downloaded from InterPro~\cite{mistryPfamProteinFamilies2021} and cleaned by removing columns with $>50\%$ gaps and sequences with $>30\%$ gaps. The resulting PCA dimensions ranged from $d{=}34$ (Pkinase) to $d{=}186$ (WW).

\begin{table}[t]
\centering
\caption{Properties of the seven Pfam families used in this study. $K$: number of seed sequences after cleaning. $L$: alignment length (columns retained). $d$: PCA dimension (95\% variance). $\bar{H}_{\mathrm{col}}$: mean per-column Shannon entropy (nats). $K_{\mathrm{eff}}$: effective number of sequences at 80\% identity threshold. $\beta^*$: empirical entropy inflection point. $\beta^*/\!\sqrt{d}$: ratio to the random-pattern prediction.}
\label{tab:families}
\small
\begin{tabular}{@{}lrrrrrrrr@{}}
\toprule
Family & Pfam ID & $K$ & $L$ & $d$ & $\bar{H}_{\mathrm{col}}$ & $K_{\mathrm{eff}}$ & $\beta^*$ & $\beta^*/\!\sqrt{d}$ \\
\midrule
RRM     & PF00076 &  68 &  71 &  59 & 1.92 &  68 & 3.85 & 0.50 \\
SH3     & PF00018 &  55 &  48 &  46 & 1.68 &  55 & 3.85 & 0.57 \\
WW      & PF00397 & 420 &  31 & 186 & 1.74 & 419 & 5.45 & 0.40 \\
Kunitz  & PF00014 &  99 &  53 &  80 & 1.61 &  99 & 3.85 & 0.43 \\
zf-C2H2 & PF00096 & 151 &  23 & 106 & 1.91 & 151 & 4.58 & 0.44 \\
PDZ     & PF00595 &  44 &  83 &  37 & 1.88 &  43 & 3.23 & 0.53 \\
Pkinase & PF00069 &  37 & 262 &  34 & 1.72 &  37 & 3.23 & 0.55 \\
\bottomrule
\end{tabular}
\end{table}

\paragraph{Scaling study.}
To characterize how generation quality degrades with decreasing family size, we fixed the WW domain (the largest family, $K{=}420$) and subsampled to $K \in \{20, 50, 100, 200, 400\}$. At each value of $K$, we drew five independent random subsets, rebuilt the memory matrix from scratch (re-computing PCA and $\beta^*$), ran stochastic attention and all baselines, and recorded metrics. This design isolates the effect of family size from family identity, and the five repeats provide error estimates.

\paragraph{Critical temperature prediction.}
We pooled $n{=}32$ data points---the 7 Pfam families and 25 WW scaling replicates (5 family sizes $\times$ 5 repeats)---and fit ordinary least-squares regression models predicting $\beta^*$ from MSA-derived features. We considered four candidate features: $\sqrt{d}$ (PCA dimensionality at 95\% retained variance), mean per-column Shannon entropy $\bar{H}_{\mathrm{col}}$, $\log K_{\mathrm{eff}}$ (effective number of sequences at 80\% identity), and spectral concentration $\lambda_1 / \mathrm{tr}(\mathbf{C})$ (leading eigenvalue of the sequence covariance, normalized by trace). We evaluated nested models by adding features in order of marginal $R^2$ improvement. Uncertainty on regression coefficients was quantified via nonparametric bootstrap: we resampled the 32 observations with replacement $10{,}000$ times, refitting the model on each resample, and reported the standard deviation of the bootstrap distribution as the coefficient SE. Model comparison used root-mean-square error (RMSE) and $R^2$, with the naive random-pattern prediction $\beta^* = \sqrt{d}$ as a reference baseline.
